% 直流潮流近似最优潮流（DC OPF）
% 本章文件由 \input 引入，不含导言区。
\section{直流潮流近似最优潮流（DC OPF）}\label{sec:dcopf}

DC OPF 是平台 OPF 家族中的线性化快速通道：把交流网络按直流潮流近似
（无损、无电压幅值变量、无无功方程）化为线性规划（LP）或凸二次规划（QP），
用于经济调度、节点边际电价（LMP）估计、可靠性评估中的大规模故障状态筛选等
对速度敏感、对电压/无功精度不敏感的场景。求解器本体为
\srcpath{src/optimal\_power\_flow/dc\_opf.cpp}（1394 行），公共入口声明于
\srcpath{include/hacdcpf/optimal\_power\_flow/dc\_opf\_solver.hpp}。与 AC OPF
（第\ref{sec:acopfnative}章、第\ref{sec:parityipm}章）不同，本求解器只覆盖
交流侧有功--相角子问题：直流网络、VSC/DC-DC/LCC 换流器物理一律不建模
（详见本章\ref{subsec:dcopf-dcgrid}节）。

\subsection{功能定位与公共接口}

\compmeta{solve\_dc\_opf / check\_dc\_opf\_feasibility}{\srcpath{include/hacdcpf/optimal\_power\_flow/dc\_opf\_solver.hpp}}

\subsubsection*{入口函数}
公共接口仅两个函数
（\srcpath{include/hacdcpf/optimal\_power\_flow/dc\_opf\_solver.hpp:solve\_dc\_opf、check\_dc\_opf\_feasibility}）：
\begin{align*}
&\text{\srcpath{solve\_dc\_opf}(\fld{sys}, \fld{opts})} \;\longrightarrow\;
  \text{\fld{DCOPFResult}}\\
&\text{\srcpath{check\_dc\_opf\_feasibility}(\fld{sys}, \fld{result}, \fld{tol})}
  \;\longrightarrow\; (\text{feasible},\ \text{max\_violation},\ \text{message})
\end{align*}
前者组装并求解 DC OPF；后者对一份已求解结果做独立的可行性复查
（发电机界、支路界、含松弛节点在内的逐节点功率平衡），容差 \fld{tol} 默认
$10^{-6}$，判据按 $\fld{tol}\times S_{\mathrm{base}}$ 折算为 MW
（\srcpath{src/optimal\_power\_flow/dc\_opf.cpp:check\_dc\_opf\_feasibility}，函数
\srcpath{check\_dc\_opf\_feasibility}；判据在 :check\_dc\_opf\_feasibility）。公共门面
\srcpath{src/api/hacdcpf.cpp:solve\_dc\_opf} 仅是透传封装，不附加语义。

\subsubsection*{内部装配结构}
建模产物集中在匿名命名空间的结构体 \fld{DCOPFFormulation}
（\srcpath{src/optimal\_power\_flow/dc\_opf.cpp:DCOPFFormulation}）：成员 \fld{lp} 与
\fld{qp} 分别承载线性化成本 LP 模型与真二次成本 QP 模型（约束结构共用），
\fld{gen\_map}/\fld{branch\_map} 记录活跃发电机/支路到原始组件向量的映射，
\fld{bus\_map} 记录母线稳定编号（\fld{ACBus::index}）到 0-based 位置的映射
（构建于 :build\_bus\_map，函数 \srcpath{build\_bus\_map}；发现重复母线编号时抛出
\fld{std::invalid\_argument}，:build\_bus\_map）。基准功率取
$\fld{base\_mva}=\max(\fld{sys.ac.base\_mva},\,1.0)$，默认 100~MVA
（:DCOPFFormulation、:build\_dc\_opf\_lp）。全章未注明处功率变量均为 $S_{\mathrm{base}}$ 标幺值。

\subsection{DC 线性化假设}\label{subsec:dcopf-assumptions}

代码中的近似假设以注释与结果声明为准，逐条如下。

\begin{enumerate}
  \item \textbf{忽略电阻、以电纳 $b=1/x$ 传输}。建模注释原文：
        ``In DC power flow: $P_{f,ij}=(\theta_i-\theta_j)/x_{ij}$ (ignore r,
        assume $x \gg r$)''（\srcpath{src/optimal\_power\_flow/dc\_opf.cpp:build\_dc\_opf\_lp}，
        函数 \srcpath{build\_dc\_opf\_lp} 的注释块 :build\_dc\_opf\_lp）。支路模型只读取
        \fld{Branch::x\_pu}，不读取电阻、分接头变比、移相角与充电电纳
        （:build\_dc\_opf\_lp）。
  \item \textbf{无电压幅值变量}。决策向量 $x=[\theta;\ P_g;\ P_f;\ \Delta P_d;\
        \lambda]$ 中只有相角 $\theta$（变量布局注释 :DCOPFFormulation），结果结构
        \fld{DCOPFResult} 也只有 \fld{va} 没有 \fld{vm}
        （\srcpath{include/hacdcpf/optimal\_power\_flow/opf\_result.hpp:DCOPFResult}），
        即经典 DC 近似中 $V_m\equiv 1$~pu 的内涵在代码里体现为“根本不存在
        电压幅值未知量”。
  \item \textbf{无无功方程}。负荷侧仅聚合有功 \fld{pd\_mw}/\fld{p\_mw}
        （:build\_dc\_opf\_lp），结果中没有任何无功字段；\fld{Generator::cost\_*} 之外的
        \fld{qmin\_mvar}/\fld{qmax\_mvar} 不进入模型。
  \item \textbf{无损}。结果对象固定写入模型局限声明：``DC OPF uses a lossless
        linearized AC network and does not model AC voltage magnitude, reactive
        power, or converter physics.''（:extract\_dc\_opf\_result）。
  \item \textbf{小角度假设未显式检验}。非松弛节点相角仅受宽界
        $-\pi\le\theta_i\le\pi$ 约束（:build\_dc\_opf\_lp），代码不做
        $|\theta_i-\theta_j|$ 越限告警；小角度是 DC 近似的理论前提而非运行期
        校验项，角度张角大的工况误差由使用者自担。
\end{enumerate}

\subsection{数学模型}\label{subsec:dcopf-model}

\subsubsection*{决策变量与布局}
设活跃母线 $n_b$、活跃发电机 $n_g$、活跃支路 $n_l$（均按 \fld{in\_service}
过滤后计数，:build\_dc\_opf\_lp），决策向量分段布局为
（:DCOPFFormulation 注释，偏移量赋值 :build\_dc\_opf\_lp）
\begin{align*}
x=[\ \underbrace{\theta}_{n_b}\ ;\ \underbrace{P_g}_{n_g}\ ;\
\underbrace{P_f}_{n_l\ \text{（可选）}}\ ;\
\underbrace{\Delta P_d}_{n_b\ \text{（可选）}}\ ;\
\underbrace{P^{gc}}_{n_b\ \text{（可选）}}\ ;\
\underbrace{\lambda}_{\text{PWL 断点}}\ ]
\end{align*}
（\srcpath{src/optimal\_power\_flow/dc\_opf.cpp:build\_dc\_opf\_lp}，函数
\srcpath{build\_dc\_opf\_lp}）。其中 $P_f$ 变量仅当
\fld{include\_branch\_limits}=\texttt{true} 时引入（:build\_dc\_opf\_lp），$\Delta P_d$
削减变量仅当 \fld{load\_shedding}=\texttt{true} 时每母线一个（:build\_dc\_opf\_lp），$P^{gc}$
仅当 \fld{allow\_fixed\_generation\_curtailment}=\texttt{true} 时每母线一个，
$\lambda$ 为凸二次成本发电机（$\fld{cost\_c2}>10^{-12}$ 且
$\fld{pmax\_mw}>\fld{pmin\_mw}$，:build\_dc\_opf\_lp）的分段线性凸组合权重，段数取
$\mathrm{clamp}(\fld{pwl\_segments},1,1000)$（:build\_dc\_opf\_lp）。

\subsubsection*{目标函数}
目标为发电成本最小化，存在 LP 与 QP 两种口径，由后端分派决定实际生效者
（见\ref{subsec:dcopf-solve}节）。

\textbf{（a）LP 口径。} 无线性化段的发电机按边际成本计价；成本数据全零时
取默认边际成本 1（成本单位/MWh），避免目标恒零
（\srcpath{src/optimal\_power\_flow/dc\_opf.cpp:build\_dc\_opf\_lp}）。由于 $P_g$ 以
标幺值进入 LP，成本系数乘以 $S_{\mathrm{base}}$ 使报告目标值与等式行对偶
保持 MW 工程单位（注释 :build\_dc\_opf\_lp）：
\begin{align*}
\min\ \sum_{k\in\mathcal{G}_{\mathrm{lin}}} c_{1,k}\,S_{\mathrm{base}}\,P_{g,k}
\;+\;\sum_{i} \mathrm{VOLL}\cdot S_{\mathrm{base}}\,\Delta P_{d,i}
\;+\;\sum_i\gamma_i S_{\mathrm{base}}P^{gc}_i
\;+\;\sum_{k\in\mathcal{G}_{\mathrm{pwl}}}\sum_{p}
f_k(\hat P_{k,p})\,\lambda_{k,p}
\end{align*}
（:build\_dc\_opf\_lp）。凸二次成本的发电机以 $\lambda$ 凸组合
插值逼近：断点
\begin{align*}
\hat P_{k,p} &= \fld{pmin\_mw}_k + \alpha_p\bigl(\fld{pmax\_mw}_k-\fld{pmin\_mw}_k\bigr),
& \alpha_p &= \frac{p}{N_{\mathrm{seg}}},\quad p=0,\dots,N_{\mathrm{seg}}\\
f_k(\hat P_{k,p}) &= c_{2,k}\hat P_{k,p}^{2}+c_{1,k}\hat P_{k,p}+c_{0,k}
\end{align*}
（:build\_dc\_opf\_lp），$\lambda_{k,p}\in[0,1]$（:build\_dc\_opf\_lp）。注释说明：凸性加
最小化使下包络自动选取相邻断点，无需二进制或 SOS2 变量（:build\_dc\_opf\_lp）。

\textbf{（b）QP 口径。} 真二次成本写成标准形
$\min \tfrac12 x^{\top}Qx+c^{\top}x$（\srcpath{src/optimal\_power\_flow/dc\_opf.cpp:build\_dc\_opf\_qp}，
函数 \srcpath{build\_dc\_opf\_qp}）：
\begin{align*}
Q_{P_{g,k},P_{g,k}} &= 2\,c_{2,k}\,S_{\mathrm{base}}^{2},
& c_{P_{g,k}} &= c_{1,k}\,S_{\mathrm{base}}
\end{align*}
（标幺换算推导见注释 :build\_dc\_opf\_qp；赋值 :build\_dc\_opf\_qp）。$c_2$、$c_1$ 同时为零时
同样回退默认边际成本 1（:build\_dc\_opf\_qp）；削减惩罚从 LP 模型原样拷贝（:build\_dc\_opf\_qp）。
常数项 $c_0$ 不进 QP 矩阵，仅在报告目标值时补加
（:compute\_qp\_objective；$c_0$ 累加亦在 :compute\_qp\_objective）。

\subsubsection*{节点功率平衡}
所有母线（含松弛节点）各有一行平衡方程
（\srcpath{src/optimal\_power\_flow/dc\_opf.cpp:build\_dc\_opf\_lp}）。显式 $P_f$ 口径下
\begin{align*}
\sum_{g\in\mathcal{G}(i)} P_{g,g} + \Delta P_{d,i}-P^{gc}_i
- \sum_{l:\,\mathrm{from}(l)=i} P_{f,l} + \sum_{l:\,\mathrm{to}(l)=i} P_{f,l}
= P^{\rm gross}_{d,i}-P^{\rm fix}_{g,i}
\end{align*}
（发电机系数 $+1$：:build\_dc\_opf\_lp；削减变量系数 $+1$ 的符号推导注释 :build\_dc\_opf\_lp；
$P_f$ 在 from 侧取 $-1$、to 侧取 $+1$：:build\_dc\_opf\_lp；右端项 $P_d$：:build\_dc\_opf\_lp）。
实现先分别聚合毛需求 $P^{\rm gross}_{d,i}$ 与固定正注入 $P^{\rm fix}_{g,i}$，再形成净右端项。
母线 \fld{pd\_mw} 和 \fld{Load::p\_mw}$\times$\fld{scaling} 的正值进入毛需求，负值进入固定注入；
\fld{StaticGenerator}、\fld{RenewableGen}、\fld{PVSystem} 与 \fld{Storage::p\_mw} 的正值进入固定注入，
负值进入毛需求。因而 $0\le\Delta P_{d,i}\le P^{\rm gross}_{d,i}$，启用 $P^{gc}$ 时
$0\le P^{gc}_i\le P^{\rm fix}_{g,i}$，不会用净负荷错误压缩可削负荷上界。

无显式 $P_f$ 变量时（\fld{include\_branch\_limits}=\texttt{false}），支路潮流
直接代入平衡行，形成经典 $B\theta$ 形式（:build\_dc\_opf\_lp）：
\begin{align*}
\sum_{j} B_{ij}\,\theta_j &= P_{g,i}+\Delta P_{d,i}-P_{d,i},
& B_{ii} &= \textstyle\sum_{j} b_{ij},\quad B_{ij}=-b_{ij},\quad b_{ij}=1/x_{ij}
\end{align*}
（$B$ 矩阵装配 :build\_dc\_opf\_lp；写入等式行 :build\_dc\_opf\_lp）。该路径曾漏掉整个松弛节点
平衡行，现仅跳过 $\theta$ 列中松弛节点的零值项而保留其平衡行（BUG-3 修复
注释 :build\_dc\_opf\_lp；对角元 :build\_dc\_opf\_lp；非对角元 :build\_dc\_opf\_lp）。

\subsubsection*{支路功率定义与输电界}
每条活跃支路一行潮流定义方程（\srcpath{src/optimal\_power\_flow/dc\_opf.cpp:build\_dc\_opf\_lp}）：
\begin{align*}
P_{f,k} - b_k\,(\theta_{\mathrm{from}(k)}-\theta_{\mathrm{to}(k)}) = 0,
\qquad b_k = 1/x_k
\end{align*}
其中 $x_k$ 取 \fld{Branch::x\_pu}，$|x_k|<10^{-12}$ 时兜底为 $10^{-6}$ 防除零
（:build\_dc\_opf\_lp；$B\theta$ 路径同兜底 :build\_dc\_opf\_lp；结果回算同兜底 :extract\_dc\_opf\_result）。
输电能力以变量盒界实现而非显式不等式（:build\_dc\_opf\_lp 注释）：
\begin{align*}
|P_{f,k}| \;\le\; \fld{rate\_a\_mva}_k\cdot\fld{branch\_limit\_margin}\,/\,S_{\mathrm{base}}
\end{align*}
（:build\_dc\_opf\_lp；\fld{rate\_a\_mva}$<10^{-9}$ 视为不限，界放宽至
\fld{kHugeBound}$=10^{20}$，:build\_dc\_opf\_lp、:kHugeBound）。注意仅长期额定 \fld{rate\_a\_mva}
生效，\fld{rate\_b}/\fld{rate\_c} 不进入模型。

\subsubsection*{发电机出力界与相角参考}
发电机界直接折算标幺盒界（:build\_dc\_opf\_lp）：
\begin{align*}
\fld{pmin\_mw}_k/S_{\mathrm{base}} \;\le\; P_{g,k} \;\le\;
\fld{pmax\_mw}_k/S_{\mathrm{base}}
\end{align*}
并以 $\pm 10^{20}$ 截断防无界（:build\_dc\_opf\_lp）。相角参考按连通分量逐分量锚定：
\srcpath{find\_slack\_bus} 先选出第一台 \fld{BusType::SLACK} 型母线，没有则
退化为第 0 号母线（:find\_slack\_bus）；再以在运支路做并查集连通分量划分，含松弛母线
的分量以松弛母线为参考，其余每个分量各自以分量内首母线为参考（:root（lambda））。
各参考相角以盒界 $[0,0]$ 钉死，其余节点自由于 $[-\pi,\pi]$（:build\_dc\_opf\_lp）。
$B\theta$ 路径保留参考列于等式约束中（其盒界已将其钉零，:build\_dc\_opf\_lp）。

\subsubsection*{负荷削减变量与 VOLL}
开启 \fld{load\_shedding} 时每母线一个削减变量
（\srcpath{src/optimal\_power\_flow/dc\_opf.cpp:build\_dc\_opf\_lp}）：
\begin{align*}
0 \;\le\; \Delta P_{d,i} \;\le\; \max\!\bigl(P_{d,i},\ 0.01/S_{\mathrm{base}}\bigr)
\end{align*}
（:build\_dc\_opf\_lp；只能削减正需求）。失负荷价值 \fld{voll}$\le 0$ 时自动取
（:build\_dc\_opf\_lp）
\begin{align*}
\mathrm{VOLL}_{\mathrm{eff}} = \max\Bigl(100,\;
1.1\max_{k}\bigl(c_{1,k}+2c_{2,k}\,\fld{pmax\_mw}_k\bigr)\Bigr)
\end{align*}
即略高于最贵发电机的边际成本上限，保证“能发不削”的字典序。

\subsubsection*{外部电网升格与固定注入}
在求解入口，每台在役 \fld{ExternalGrid} 被升格为一台挂在原母线上的松弛发电机，
界 $\pm 10^{5}$~MW/MVAr，成本系数照搬 \fld{ExternalGrid::cost\_c2/c1/c0}
（\srcpath{src/optimal\_power\_flow/dc\_opf.cpp:solve\_dc\_opf}，函数
\srcpath{solve\_dc\_opf}）；求解后其出力拆回
\fld{external\_grid\_p\_mw}，并从 \fld{pg\_mw} 中裁掉（:separate\_external\_grid\_dispatch（lambda），定义于
\srcpath{solve\_dc\_opf}）。

\subsection{求解流程与后端分派}\label{subsec:dcopf-solve}

\subsubsection*{求解主流程}
\srcpath{solve\_dc\_opf}（\srcpath{src/optimal\_power\_flow/dc\_opf.cpp:solve\_dc\_opf}）
的总体流程如图~\ref{fig:dcopf-flow}：先做 LCC 拒绝与规范投影，再做图论孤岛
预检，然后一次性装配 LP 约束结构并派生 QP 模型，最后按后端枚举分派求解并
回收对偶。

\begin{figure}[H]
\centering
\begin{tikzpicture}[
  node distance=0.55cm,
  box/.style={draw, rectangle, align=center, inner sep=3.5pt, font=\small},
  dec/.style={draw, diamond, aspect=2.1, align=center, inner sep=1.5pt, font=\small},
  arr/.style={-{Stealth[length=2.2mm]}, thick}]
\node[box] (in) {输入 \fld{HybridPowerSystem}};
\node[dec, below=of in] (lcc) {含在役 LCC？};
\node[box, right=1.6cm of lcc] (rej) {拒绝求解\\ \fld{dc-opf:rejected-lcc}};
\node[box, below=of lcc] (proj) {规范投影（开关合并母线）\\ + 外网升格为松弛发电机};
\node[dec, below=of proj] (isl) {孤岛预检\\ 存在无松弛孤岛？};
\node[box, right=1.6cm of isl] (shed) {预削减孤岛负荷\\ 无有效孤岛则报不可行};
\node[box, below=of isl] (build) {装配 LP 约束结构\\ 派生 QP 成本模型};
\node[box, below=of build] (disp) {后端分派（回退链）\\ Gurobi / LCQP / HiGHS / 对偶单纯形};
\node[box, below=of disp] (dual) {支撑 LP 回收对偶（LMP）\\ \fld{compute\_lmp} 门控};
\node[box, below=of dual] (out) {结果归因：反投影回原始母线/支路\\ 填写 \fld{model\_limitations} 等诚实字段};
\draw[arr] (in) -- (lcc);
\draw[arr] (lcc) -- node[above, font=\small]{是} (rej);
\draw[arr] (lcc) -- node[right, font=\small]{否} (proj);
\draw[arr] (proj) -- (isl);
\draw[arr] (isl) -- node[above, font=\small]{是} (shed);
\draw[arr] (isl) -- node[right, font=\small]{否/剪枝后} (build);
\draw[arr] (shed.south) |- (build.east);
\draw[arr] (build) -- (disp);
\draw[arr] (disp) -- (dual);
\draw[arr] (dual) -- (out);
\end{tikzpicture}
\caption{DC OPF 求解主流程（对应
\srcpath{src/optimal\_power\_flow/dc\_opf.cpp:solve\_dc\_opf}）}
\label{fig:dcopf-flow}
\end{figure}

\textbf{LCC 拒绝}：任何在役 \fld{LCCConverter} 都使求解直接返回，状态写明
“this AC-only linear formulation does not model LCC AC/DC coupling”，并提示改用
ParityIPM 或 Ipopt 的 AC OPF（:solve\_dc\_opf；见第\ref{sec:parityipm}章）。

\textbf{规范投影}：与 AC OPF/潮流共用同一份 canonical 拓扑——闭合开关与
断路器合并母线（\fld{strip\_dead\_islands}=\texttt{false}，:solve\_dc\_opf）。求解后
\fld{va}/\fld{lmp} 按密集量、\fld{load\_shedding\_mw} 按广延量反投影回原始
母线，\fld{pf\_mw} 按 \fld{branch\_orig\_to\_proj} 映射回原始支路（:separate\_external\_grid\_dispatch（lambda））。
失败解不回灌投影，以免用尺寸异常掩盖原始不可行性（注释 :separate\_external\_grid\_dispatch（lambda））。

\textbf{孤岛预检}：用图模块 \srcpath{build\_power\_system\_graph} 与
\srcpath{analyze\_topology} 在昂贵的 LP 装配前识别死岛——\fld{IsolatedLoad}
孤岛，或 \fld{NoSlack} 且岛内无在役可调度发电机的孤岛（:solve\_dc\_opf）；
\fld{NoSlack} 但含可调度电源的孤岛判活并正常建模。死岛负荷（母线 \fld{pd\_mw} 与
\fld{ac.loads} 两源相加，口径与建模一致，P1a 注释 :solve\_dc\_opf）计入直接削减
（:solve\_dc\_opf）；全部孤岛均无效时直接返回
\texttt{"Infeasible: no island with slack bus"}（:solve\_dc\_opf）；否则把死岛上的
负荷、发电机、静态发电机、新能源、光伏、储能与相连支路在本地副本中置为
退出再建模（:solve\_dc\_opf）。

\textbf{装配}：\srcpath{build\_dc\_opf\_lp} 构建约束结构（:solve\_dc\_opf 调用；
重复母线编号等拓扑错误在此捕获并转状态，:solve\_dc\_opf）；空系统或无发电机
直接返回 \texttt{"Empty system or no generators"}（:solve\_dc\_opf）；随后
\srcpath{build\_dc\_opf\_qp} 派生真二次成本模型（:solve\_dc\_opf）。

\subsubsection*{后端选择与回退链}
后端枚举 \fld{DCOPFSolverBackend} 取值为 \fld{Auto}/\fld{Native}/
\fld{NativeQP}/\fld{HiGHS}/\fld{Gurobi}
（\srcpath{include/hacdcpf/optimal\_power\_flow/opf\_options.hpp:DCOPFSolverBackend}）。
各候选通道（均定义于 \srcpath{src/optimal\_power\_flow/dc\_opf.cpp}）：
\begin{itemize}
  \item \textbf{Gurobi QP}（:try\_gurobi\_qp（lambda），定义于 \srcpath{solve\_dc\_opf}）：
        经 engine 适配器 \fld{GurobiAdapter} 求解 QP；仅在
        \fld{gurobi.available()} 探测到安装/许可证时尝试，否则记
        \fld{unavailable}；
  \item \textbf{NativeLCQP}（:try\_native\_qp（lambda），定义于 \srcpath{solve\_dc\_opf}）：
        自研线性约束 QP 内点求解器（MIPSolvers engine 的
        \srcpath{lcqp\_solver}，经转发头
        \srcpath{include/hacdcpf/engine/kernel/ipm/lcqp\_solver.hpp} 引入），
        求解真二次成本 QP；
  \item \textbf{HiGHS 通道}（:try\_highs（lambda），定义于 \srcpath{solve\_dc\_opf}）：
        HiGHS 本身无原生 QP 接口，先经 \fld{SolverEngine} 以
        \fld{preferred\_solver}=\fld{"HiGHS"} 路由一次 QP（实际求解者可能是
        Gurobi 或 LCQP，\fld{solver\_name} 如实反映，注释 :try\_highs（lambda））；QP
        失败后回退 \fld{HighsAdapter::solve\_lp} 求 LP，此时目标与对偶只反映
        线性成本模型，求解器名记为 \fld{"HiGHS-LP(QP-fallback)"}
        （:try\_highs（lambda））；
  \item \textbf{NativeDualSimplex}（:try\_native\_simplex（lambda），定义于
        \srcpath{solve\_dc\_opf}）：自研对偶单纯形
        \srcpath{solve\_lp\_with\_basis}（同样来自 MIPSolvers engine，经转发头
        \srcpath{include/hacdcpf/engine/kernel/lp\_kernel/dual\_simplex.hpp}），
        求解线性化成本 LP。
\end{itemize}
回退顺序（:solve\_dc\_opf）：\fld{Auto} 依次试 Gurobi-QP → NativeLCQP → HiGHS
通道 → NativeDualSimplex（:solve\_dc\_opf）；\fld{Gurobi} 同序但从 Gurobi 起
（:solve\_dc\_opf）；\fld{NativeQP} 失败仅回退 NativeDualSimplex（:solve\_dc\_opf）；
\fld{HiGHS} 不可用时直接失败返回（:solve\_dc\_opf）；\fld{Native} 只用对偶单纯形
（:solve\_dc\_opf）。每一步尝试都记入 \fld{solver\_chain}，格式
\fld{<backend>:<ok|fail(detail)>}，供事后审计真实回退路径（:record\_chain（lambda）、
:solve\_dc\_opf）。

\textbf{目标语义消歧}：LP 回退成功时必须把 \fld{use\_qp} 复位，否则会用 QP
口径重算 LP 可行点，造成目标/对偶语义不一致（注释 :solve\_dc\_opf）。最终按收敛路径重算或保留目标值并填写
\fld{objective\_model}：QP 收敛为 \fld{"QP"}（真二次成本含 $c_0$，由
\srcpath{compute\_qp\_objective} 重算，:solve\_dc\_opf），否则为 \fld{"LP-PWL"}
或 \fld{"LP"}，并以 \fld{pwl\_segments\_effective} 记录实际分段数（QP 时置 0，
:solve\_dc\_opf）。

\subsubsection*{对偶恢复与支撑 LP}
QP/LP 主解的对偶若不完整（\srcpath{has\_complete\_duals} 校验行对偶与盒对偶
数量，:has\_complete\_duals），且 \fld{compute\_lmp}=\texttt{true} 时，追加一次支撑 LP：
以原 LP 模型（QP 情形把成本换成梯度 $Qx+c$）用对偶单纯形重解，取回全部行
对偶与盒对偶（:populate\_missing\_duals\_from\_supporting\_lp；QP 梯度替换
:populate\_missing\_duals\_from\_supporting\_lp；门控 :solve\_dc\_opf）。代码注释如实声明该通道的已知缺陷：自研单纯形对
有界变量 LP 的盒对偶恢复不正确，\fld{include\_branch\_limits}=\texttt{true}
且支路界真正起作用时 \fld{branch\_mu\_lower}/\fld{branch\_mu\_upper}
\textbf{不可靠}；LMP（节点平衡行对偶）仍有效可用于经济调度
（NOTE 注释块 :populate\_missing\_duals\_from\_supporting\_lp）。因此结果中 \fld{branch\_mu\_valid} 在 KKT 级提取
实现之前\textbf{恒为 false}（:extract\_dc\_opf\_result）。

\subsubsection*{与 HiGHS/AML 建模层的分派关系}
本求解器只经 engine 适配器调用 HiGHS/Gurobi（转发头
\srcpath{include/hacdcpf/engine/solver/external/adapters.hpp}），自身不经过
AML 建模层。平台另有一套基于 MIPSolvers AML 的 DC OPF 实现
\srcpath{power\_models::solve\_dc\_opf}
（\srcpath{src/power\_models/dc\_opf\_builder.cpp:solve\_dc\_opf}），以集合/参数/变量的
代数建模语言重构同一数学模型并交由 AML 后端求解，用于交叉验证；其建模细节
仅见独立的 \srcpath{docs/modules/power\_models/power\_models\_manual.tex}。
两套实现相互独立：本章所述的孤岛预检、外网
升格、负荷削减与支撑 LP 对偶恢复只存在于引擎版。

\subsection{直流电网与换流器覆盖情况}\label{subsec:dcopf-dcgrid}

\textbf{不覆盖。} 具体口径：
\begin{itemize}
  \item 直流母线、直流支路、VSC 换流器、DC/DC 变换器、能量路由器均不进入
        模型——建模仅遍历 \fld{sys.ac.*} 容器（
        \srcpath{src/optimal\_power\_flow/dc\_opf.cpp:build\_dc\_opf\_lp}），结果结构也
        无任何直流侧字段
        （\srcpath{include/hacdcpf/optimal\_power\_flow/opf\_result.hpp:DCOPFResult}）；
  \item 结果固定声明 \fld{model\_scope}=\fld{"dc-opf:linear-no-converter-model"}
        （:extract\_dc\_opf\_result），\fld{converter\_model\_scope} 各换流器标志位恒为
        \texttt{false}（注释见
        \srcpath{include/hacdcpf/optimal\_power\_flow/opf\_result.hpp:DCOPFResult}）；
  \item 含在役 LCC 的系统被整体拒绝而非近似处理（:solve\_dc\_opf）；
  \item 直流侧负荷由此形成已知盲区：可靠性评估调用方对此有显式告警
        （“DC 负荷中断不计入 OPF，EENS/LOLE 可能偏乐观”，
        \srcpath{src/reliability/reliability\_assessment.cpp:evaluate\_state}）。
\end{itemize}
需要交直流联合优化的场景应使用 Parity IPM 全空间 OPF
（第\ref{sec:parityipm}章）或 AML 层的 acdcopf builder；后者见独立的
\srcpath{docs/modules/power\_models/power\_models\_manual.tex}。

\subsection{NativeLCQP 的结构化初值}

对每个在役 AC 支路连通分量 $\mathcal C$，先在发电上下界内调整 $P_g$，不足
部分再由允许的 $\Delta P_d$ 补足，使
\begin{equation}
\sum_{i\in\mathcal C}(P_{g,i}+\Delta P_{d,i}-P_{d,i})=0.
\end{equation}
每个分量选择一个参考角并删除其行列，求约化 Laplacian
$B_{\mathcal C,r}\theta_r=p_{\mathcal C,r}$；随后由
$P_{f,ij}=(\theta_i-\theta_j)/x_{ij}$ 回填显式支路流，并补齐 PWL convexity/
interpolation 变量。因此连通分量功率和是 Laplacian 可解性的兼容条件，而非
额外优化问题。生成点经过变量盒界与 $\|A_{eq}x_0-b_{eq}\|_\infty$ 审计后，
只通过 \fld{QPModel::x0} 交给 NativeLCQP；其他后端不消费该初值，构建失败则
精确回到确定性冷初始化。该步骤复杂度为图遍历线性工作加各分量一次约化
Laplacian 稀疏分解，不运行非线性 Phase I。

作为 ACOPF warm-start-only 子问题时，
\fld{phase\_one\_time\_limit\_ms} 从 \srcpath{solve\_dc\_opf} 入口开始计时，
因此 canonical projection、formulation、上述结构初值、symbolic analyze 与
numeric iteration 都计入总预算。NativeLCQP 在迭代/分解边界协作式检查截止；
一个已经开始的稀疏分解不可中断，允许完成后报告
\fld{phase\_one\_budget\_overshoot\_ms}。超时迭代仍须通过显式原坐标残差门槛
才可作为 warm start，且不声明 DCOPF 最优；有限预算耗尽后不再运行 simplex
fallback。

\subsection{选项：DCOPFOptions}

\compmeta{DCOPFOptions}{\srcpath{include/hacdcpf/optimal\_power\_flow/opf\_options.hpp}}

全部字段如下（定义 :DCOPFOptions）。

\begin{paramtable}
\fld{feasibility\_tol} & $\varepsilon_{\mathrm{f}}$ & — & $10^{-9}$--$10^{-4}$（经验值） & $10^{-6}$ & 原始可行性容差；传入 LCQP（\fld{tol\_primal}）与对偶单纯形\\
\fld{optimality\_tol} & $\varepsilon_{\mathrm{o}}$ & — & $10^{-9}$--$10^{-4}$（经验值） & $10^{-6}$ & 对偶/最优性容差；传入 LCQP（\fld{tol\_dual}/\fld{tol\_gap}）与对偶单纯形\\
\fld{max\_iterations} & — & — & $10^{2}$--$10^{5}$（经验值） & 10000 & 单纯形/LCQP 最大迭代数\\
\fld{verbose} & — & — & true/false & false & spdlog 打印规模与收敛信息（:solve\_dc\_opf）\\
\fld{solver} & — & — & Auto/Native/\allowbreak NativeQP/\allowbreak HiGHS/\allowbreak Gurobi & Auto & 后端选择与回退链入口（:DCOPFOptions）\\
\fld{pwl\_segments} & $N_{\mathrm{seg}}$ & — & 1--1000（clamp，:build\_dc\_opf\_lp） & 4 & 凸二次成本的 LP 分段线性插值段数；QP 路径不使用\\
\fld{include\_branch\_limits} & — & — & true/false & true & 是否引入显式 $P_f$ 变量及 \fld{rate\_a\_mva} 盒界；false 时退回纯 $B\theta$ 形式且无支路约束\\
\fld{branch\_limit\_margin} & — & — & 0.9--1.0（经验值） & 1.0 & 支路界缩放系数，乘在 \fld{rate\_a\_mva} 上（:DCOPFOptions）\\
\fld{load\_shedding} & — & — & true/false & true & 是否引入每母线削减变量 $\Delta P_d$\\
\fld{voll} & VOLL & —/MWh & $\ge 0$ & 0.0 & 失负荷价值；$\le 0$ 时按 $\max(100,\,1.1\times$最高边际成本$)$ 自动取值（:DCOPFOptions）\\
\fld{lexicographic\_load\_shedding} & — & — & true/false & false & 两阶段精确优先级：先最小总削减，再固定最优削减量并最小发电成本；HiGHS 第二阶段使用 PWL 成本 LP\\
\fld{full\_redispatch\_from\_zero} & — & — & true/false & false & HL-I/HL-II 充裕度再调度语义，在线机组统一采用 $0\le P_g\le P_g^{\max}$\\
\fld{allow\_fixed\_generation\_curtailment} & — & — & true/false & false & 是否引入逐母线 $P^{gc}$；与全负荷削减、零机组下界共同给出任意拓扑的构造性可行点\\
\fld{fixed\_generation\_curtailment\_cost} & $\gamma$ & 成本/MWh & $\ge0$ & 1.0 & 第二级目标中的固定正注入削减代价；第一级最小切负荷不使用\\
\fld{compute\_lmp} & — & — & true/false & true & 主解后是否追加支撑 LP 回收对偶；false 时 \fld{lmp} 清空并注明原因（:solve\_dc\_opf）\\
\fld{structural\_warm\_start} & — & — & true/false & true & 为 NativeLCQP 构造逐分量功率配平与约化 Laplacian 初值；LP/外部后端忽略\\
\fld{phase\_one\_linear\_relaxation} & — & — & true/false & false & 仅 HiGHS 后端的 Phase-I 凸松弛：直接解 LP-PWL 而非把二次模型路由到 QP 后端；返回点为 warm start，非二次成本 DCOPF 最优证书\\
\fld{compact\_quadratic\_model} & — & — & true/false & false & Phase-I QP 省略仅供 LP-PWL 使用的插值变量与行\\
\fld{accept\_phase\_one\_iterate} & — & — & true/false & false & 允许非最优 NativeLCQP 末迭代经显式残差门槛成为 warm-start-only 结果\\
\fld{phase\_one\_iterate\_tolerance} & $\varepsilon_I$ & pu & $\ge0$ & 0.1 & warm-start-only 末迭代的 equality/box 最大残差门槛\\
\fld{phase\_one\_time\_limit\_ms} & $T_I$ & ms & $\le0$ 或 $>0$ & 0 & warm-start-only 端到端协作式墙钟预算；$\le0$ 不限时\\
\end{paramtable}

\subsection{结果字段：DCOPFResult}

\compmeta{DCOPFResult}{\srcpath{include/hacdcpf/optimal\_power\_flow/opf\_result.hpp}}

全部字段如下（定义 :DCOPFResult）。向量字段的索引空间：\fld{va}/\fld{lmp}/
\fld{load\_shedding\_mw}/\fld{fixed\_generation\_curtailment\_mw} 按原始母线向量位置（经反投影）；\fld{pg\_mw} 按
\fld{sys.ac.generators} 位置；\fld{pf\_mw}/\fld{branch\_mu\_*} 按
\fld{sys.ac.branches} 位置；\fld{external\_grid\_p\_mw} 按
\fld{sys.ac.external\_grids} 位置。

\begin{paramtable}
\fld{va} & $\theta_i$ & rad & $(-\pi,\pi)$ & 空 & 母线电压相角（各连通分量参考母线处恒 0）\\
\fld{pg\_mw} & $P_g$ & MW & \fld{pmin\_mw}--\fld{pmax\_mw} & 空 & 发电机有功调度（由标幺乘 $S_{\mathrm{base}}$ 还原，:extract\_dc\_opf\_result）\\
\fld{external\_grid\_p\_mw} & $P_{\mathrm{ext}}$ & MW & $\pm 10^{5}$ & 空 & 外网（升格发电机）有功（:DCOPFResult）\\
\fld{pf\_mw} & $P_f$ & MW & $\pm$\fld{rate\_a\_mva} & 空 & 支路有功，from$\to$to 为正；无显式 $P_f$ 时由 $(\theta_f-\theta_t)/x$ 回算（:DCOPFResult）\\
\fld{converged} & — & — & true/false & false & 后端成功且规范坐标原始可行性证书通过时为 true\\
\fld{iterations} & — & — & — & 0 & 最终后端迭代数\\
\fld{objective} & $f$ & —/h & — & 0.0 & 目标值；口径见 \fld{objective\_model}（QP 含 $c_0$，LP 为线性/PWL 成本）\\
\fld{status} & — & — & — & 空 & 后端状态串或拒绝/失败原因（如 \texttt{"HiGHS not available"}）\\
\fld{solver\_name} & — & — & — & 空 & 实际收敛后端名（可能含 \fld{(QP-fallback)} 后缀）\\
\fld{runtime\_sec} & — & s & — & 0.0 & 含投影与预检的端到端耗时（:DCOPFResult）\\
\fld{structural\_warm\_start\_*} & — & pu/计数/布尔 & — & inf/0/false & 请求、构建、实际消费、分量/分解数、原坐标 equality residual 与状态原因\\
\fld{solver\_initial\_primal\_residual} & — & 缩放坐标 & $\ge0$ & inf & NativeLCQP 第 0 次迭代的 primal residual；含严格正 bound-slack 残差\\
\fld{phase\_one\_warm\_start\_only} & — & — & true/false & false & 末迭代只获准供上层 warm start 使用，不是 DCOPF 最优证书\\
\fld{phase\_one\_budget\_exhausted} & — & — & true/false & false & formulation 或 LCQP 阶段达到有限墙钟预算\\
\fld{phase\_one\_budget\_overshoot\_ms} & — & ms & $\ge0$ & 0 & 不可中断在途工作造成的预算超时\\
\fld{native\_qp\_symbolic\_analyze\_calls} & — & 次 & 0--1 & 0 & 本次 NativeLCQP KKT 稀疏模式分析次数；每轮 numeric factorization 复用该模式\\
\fld{lmp} & $\lambda_i$ & —/MWh & $\ge 0$（无罚项时） & 空 & 节点边际电价：平衡行对偶除以 $S_{\mathrm{base}}$（:DCOPFResult）\\
\fld{branch\_mu\_lower} & $\mu^{-}$ & —/MWh & — & 空 & 支路下界盒对偶/$S_{\mathrm{base}}$；当前\textbf{不可靠}（:DCOPFResult）\\
\fld{branch\_mu\_upper} & $\mu^{+}$ & —/MWh & — & 空 & 支路上界盒对偶/$S_{\mathrm{base}}$；当前\textbf{不可靠}（:DCOPFResult）\\
\fld{branch\_mu\_valid} & — & — & true/false & false & 恒 false，直至 KKT 级盒对偶提取实现（:extract\_dc\_opf\_result）\\
\fld{branch\_mu\_validity\_reason} & — & — & — & 空 & 上述无效性的人类可读说明\\
\fld{lmp\_valid} & — & — & true/false & false & 行对偶证书完整时为 true（:DCOPFResult）\\
\fld{lmp\_validity\_reason} & — & — & — & 空 & LMP 来源/缺失原因说明\\
\fld{load\_shedding\_mw} & $\Delta P_d$ & MW & $\ge 0$ & 空 & 逐母线模型内生削减量；$<10^{-6}$ 数值噪声清零\\
\fld{total\_load\_shedding\_mw} & — & MW & $\ge 0$ & 0.0 & 削减总量\\
\fld{fixed\_generation\_curtailment\_mw} & $P^{gc}$ & MW & $\ge0$ & 空 & 逐母线固定正注入削减量\\
\fld{total\_fixed\_generation\_curtailment\_mw} & — & MW & $\ge0$ & 0.0 & 固定正注入削减总量\\
\fld{full\_redispatch\_from\_zero} & — & — & true/false & false & 回显本次有效机组下界语义\\
\fld{primal\_feasibility\_certified} & — & — & true/false & false & 反投影前在实际规范 LP/QP 坐标中认证全部等式、不等式和盒界\\
\fld{maximum\_primal\_violation\_mw} & — & MW & $\ge0$ & inf & 上述规范坐标证书的最大违反量\\
\fld{primal\_feasibility\_reason} & — & — & — & 空 & 证书通过或拒绝原因\\
\fld{solver\_chain} & — & — & — & 空 & 依序记录各后端尝试结果，供审计回退路径（:DCOPFResult）\\
\fld{objective\_model} & — & — & QP/LP-PWL/LP & 空 & 报告目标值对应的成本口径（:DCOPFResult）\\
\fld{pwl\_segments\_effective} & — & — & 0--1000 & 0 & LP-PWL 实际段数；0 表示 QP 真二次成本\\
\fld{converter\_model\_scope} & — & — & — & 全 false & 换流器建模覆盖声明（DC OPF 不建模任何换流器）\\
\fld{model\_limitations} & — & — & — & 空 & 本次求解的模型局限声明（无损线性化、储能单时段等，:extract\_dc\_opf\_result）\\
\end{paramtable}

\subsection{典型使用场景}

\begin{itemize}
  \item \textbf{可靠性评估（FMEA/状态评估）}：每个故障状态开启毛负荷削减、固定正注入削减和
        零机组下界；先解纯最小切负荷 LP，再加入最优切负荷等式并最小化 PWL 发电与 $P^{gc}$ 成本，
        不使用有限 VOLL 冒充字典序
        （\srcpath{src/reliability/reliability\_assessment.cpp:evaluate\_state} 调用，
        选项设置 :run\_nonsequential\_mc；AC-only 适用范围声明 :evaluate\_state）。
  \item \textbf{EV--交通耦合联合优化}：按时段把 EV 充电负荷注入系统后逐时段
        求解 DC OPF，累加发电成本并取 LMP 构成联合社会福利目标
        （\srcpath{src/ev\_power\_traffic/joint\_social\_welfare.cpp:solve\_joint\_social\_welfare}；
        CTM 变体 \srcpath{src/ev\_power\_traffic/ctm\_joint\_welfare.cpp:simulate\_ev\_power\_traffic\_ctm\_joint}）。
  \item \textbf{SPPT 语义保持验证}：MR3d 用 DC OPF 节点电价校验
        rich$\to$canonical 投影前后的语义不变性
        （\srcpath{src/sppt/metamorphic.cpp:mr3\_semantic\_preservation\_opf\_dual}）。
  \item \textbf{HTTP 后端}：v1 分析接口在 \fld{analysis}=\fld{"optimal\_power\_flow"}
        且 \fld{solver}=\fld{"dc"} 时调用本求解器并序列化结果
        （\srcpath{src/server/runtime\_api\_v1.cpp:execute\_analysis}，函数
        \srcpath{execute\_analysis}）。
  \item \textbf{公共门面}：\srcpath{src/api/hacdcpf.cpp:solve\_dc\_opf} 透传封装。
\end{itemize}
如实说明两处“不调用”：\textbf{time\_series} 多时段/年度流水线挂接的是
AC OPF（\srcpath{src/time\_series/time\_series\_pf.cpp} 引入
\srcpath{ac\_opf\_solver.hpp}），不使用本求解器；\textbf{market} 模块的
SCUC/SCED 为其自研 LP 体系（\srcpath{src/market/market\_simulation.cpp:build\_pricing\_model}
附近），亦不复用 \srcpath{solve\_dc\_opf}。

\subsection{验证与校核}

\begin{itemize}
  \item \textbf{MATPOWER 标准算例}：case9 收敛且目标有限、\fld{pg\_mw} 尺寸
        正确（\srcpath{tests/test\_matpower\_cases.cpp:TEST\_CASE("DC OPF: case9 solves and objective is finite")}）；case14 调度
        逐机检界（同文件 :TEST\_CASE("DC OPF: case14 dispatch within generator limits")）。
  \item \textbf{AC/DC/PF 三方交叉验证}：小算例集上 DC OPF 收敛后以 AC 潮流
        回代其调度检验可行性；case57 上“DC 调度$\to$AC 潮流”不保证收敛，
        仅告警不断言（\srcpath{tests/test\_acopf\_dcopf\_crossval.cpp:TEST\_CASE("ACOPF/DCOPF/PF cross-validation on small cases")}，
        豁免逻辑 :TEST\_CASE("ACOPF/DCOPF/PF cross-validation on small cases")）。
  \item \textbf{可行性自检}：\srcpath{check\_dc\_opf\_feasibility} 对 case14
        解复检通过（同文件 :TEST\_CASE("DC OPF solution feasibility")）。
  \item \textbf{孤岛回归}：孤立负荷岛逐母线削减记账（:TEST\_CASE("DC OPF isolated-island shedding is reported per bus")）、无 authored
        松弛但含可调度电源的孤岛\textbf{被保留}、岛内机组出力承担岛内负荷而不做削减
        （:TEST\_CASE("DC OPF preserves a source-capable island without an authored slack")）、松弛节点平衡行纳入复检（:TEST\_CASE("DC OPF feasibility check includes slack-bus balance")）。
  \item \textbf{调度质量}：case30 上与 AC OPF（NativeAC）对比，$\sum|\Delta
        P_g|/\sum P_g$ 容忍阈值为 1.0——注释明言“DC OPF is a crude
        approximation”（同文件 :TEST\_CASE("DC OPF dispatch quality")，阈值同该用例）。
  \item \textbf{LMP 口径}：无支路界时 LMP 非负且不超过最高边际成本的数量级
        （防 $S_{\mathrm{base}}$ 换算 bug，同文件 :TEST\_CASE("DC OPF LMP: sign and scale")）；无拥塞时全网
        LMP 一致（:TEST\_CASE("DC OPF LMP: uncongested uniformity")）。
\end{itemize}
验证体系的总体组织见第\ref{sec:verification}章。

\subsection{边界与局限}

\begin{enumerate}
  \item \textbf{DC 近似的固有盲区}：无损、无电压幅值、无无功（结果
        \fld{model\_limitations} 首条，\srcpath{src/optimal\_power\_flow/dc\_opf.cpp:extract\_dc\_opf\_result}）。
        电压越限、无功不足、网损引起的附加阻塞均不可见；重载、大角度张角或
        低 $x/r$ 比配电网场景下误差显著，代码不做任何运行期误差估计。
  \item \textbf{支路模型简化}：仅取 \fld{x\_pu}，忽略电阻、分接头、移相角与
        充电电纳（:build\_dc\_opf\_lp）；限值只用 \fld{rate\_a\_mva}（:build\_dc\_opf\_lp）；零电抗支路
        以 $x=10^{-6}$ 兜底而非合并（:build\_dc\_opf\_lp），可能引入数值
        病态。
  \item \textbf{支路拥塞对偶不可用}：\fld{branch\_mu\_valid} 恒 false
        （:extract\_dc\_opf\_result）；需要拥塞租金分解的研究只能使用 LMP 总量，不能拆分
        能量/拥塞分量（NOTE 注释 :populate\_missing\_duals\_from\_supporting\_lp）。
  \item \textbf{LMP 的适用范围}：LMP 取自支撑 LP 平衡行对偶，标幺换算已按
        $1/S_{\mathrm{base}}$ 还原（:extract\_dc\_opf\_result）；\fld{include\_branch\_limits}
        =\texttt{true} 时 LMP 数值仍有效，但注意测试约定 LMP 专项用例是在
        \fld{include\_branch\_limits}=\texttt{false} 下跑的
        （\srcpath{tests/test\_acopf\_dcopf\_crossval.cpp:TEST\_CASE("DC OPF LMP: sign and scale")} 注释）。
  \item \textbf{无直流网络与换流器}：见本章\ref{subsec:dcopf-dcgrid}节；
        直流负荷对结果完全不可见。
  \item \textbf{储能为固定注入}：\fld{Storage::p\_mw} 作为不可调注入扣除，
        无跨时段 SOC 动态（:extract\_dc\_opf\_result 如实声明）。
  \item \textbf{外网升格为宽界松弛发电机}：$\pm 10^{5}$~MW（:solve\_dc\_opf），
        多外网并存时各外网出力由成本排序决定，不模拟实际互联调度协议。
  \item \textbf{后端可用性依赖环境}：Gurobi 需本机安装与许可证
        （:try\_gurobi\_qp（lambda））；\fld{HiGHS} 后端在库不可用时直接失败
        （:solve\_dc\_opf）；QP 经 \fld{SolverEngine} 路由时实际求解者以
        \fld{solver\_name}/\fld{solver\_chain} 为准（:try\_highs（lambda））。
  \item \textbf{LP-PWL 插值误差}：凸二次成本走 LP 路径时为分段线性近似，
        精度由 \fld{pwl\_segments}（默认 4）控制；报告口径以
        \fld{objective\_model}/\fld{pwl\_segments\_effective} 区分
        （:solve\_dc\_opf）。
  \item \textbf{相角界为宽界}：$[-\pi,\pi]$ 仅防无界（:build\_dc\_opf\_lp），不构成
        静态安全（角差）约束；需要角差约束的研究应改用
        AC OPF 或 AML 建模层（见
        \srcpath{docs/modules/power\_models/power\_models\_manual.tex}）。
\end{enumerate}
